\section{Thermodynamically consistent variable-gas formulations}
\label{sec:thermo}
\subsection{The gas model is a constitutive system}
The current MOC uses constant-\(\gamma\) Prandtl--Meyer functions,
isentropic power laws, characteristic coefficients and transonic starting
data. Replacing \(\gamma\) at individual stations leaves these ingredients
mathematically inconsistent unless the equations are rederived for the
selected thermodynamic regime. Anderson explicitly distinguishes an
effective \(\gamma\) fitted to one relation from an exact representation
of all the other relations \cite[\S17.6, PDF pp.~684--685]{anderson}.

\begin{longtable}{@{}P{0.20\linewidth}P{0.35\linewidth}P{0.37\linewidth}@{}}
\toprule\textbf{Regime} & \textbf{Closure} & \textbf{MOC implication}\\\midrule
\endhead
Calorically perfect &
Fixed composition and constant \(c_p,c_v,R,\gamma\). &
Present branch; retain as the limiting-case regression.\\
Thermally perfect, frozen composition &
Fixed species mass fractions \(Y_i\), but \(c_p(T)\) varies.
A single temperature and internal-mode equilibrium are explicit choices. &
First practical variable-property branch; use a consistent EOS,
enthalpy, entropy and frozen acoustic speed.\\
Chemical equilibrium &
Composition follows local equilibrium subject to element conservation. &
Requires constrained equilibrium derivatives and the appropriate
equilibrium acoustic closure. A frozen \(c_p(T)\) table is insufficient.\\
Finite-rate nonequilibrium &
Composition and possibly internal-mode energies evolve along flow history. &
Adds species/relaxation transport and stiffness. Entropy need not be
constant in inviscid adiabatic expansion.\\
\bottomrule
\end{longtable}

These distinctions follow the original thermodynamics discussion, not
software naming conventions \cite[\S\S17.4--17.9, PDF pp.~675--689]{anderson}.
An equilibrium/frozen label must specify what is frozen, at which station,
and with which composition, internal-mode and thermodynamic data.

\subsection{First target: fixed-composition thermally perfect gas}
For a single-phase ideal mixture with fixed \(Y_i\),
\begin{align}
R &= \sum_iY_i R_i, & p&=\rho RT,\\
c_p(T)&=\sum_iY_i c_{p,i}(T),&
c_v(T)&=c_p(T)-R,\\
h(T)-h(T_{\mathrm{ref}})
&=\int_{T_{\mathrm{ref}}}^{T}c_p(\tau)\,\dd\tau,\\
s(T,p)-s(T_{\mathrm{ref}},p_{\mathrm{ref}})
&=\int_{T_{\mathrm{ref}}}^{T}\frac{c_p(\tau)}{\tau}\,\dd\tau
-R\ln\!\left(\frac{p}{p_{\mathrm{ref}}}\right).
\end{align}
The original text gives the thermally perfect entropy integral and
enthalpy identities explicitly \cite[\S1.4.5, PDF pp.~40--44,
Eqs.~1.18, 1.20, 1.32--1.35; \S17.4, PDF p.~678]{anderson}.
The frozen small-disturbance sound speed is
\begin{equation}
a_f^2=\left(\frac{\partial p}{\partial\rho}\right)_{s,Y}
=\gamma_f(T)RT,\qquad
\gamma_f=\frac{c_{p,f}}{c_{v,f}},
\end{equation}
under these assumptions \cite[\S17.5, PDF p.~682, Eq.~17.36]{anderson}.
This acoustic definition, rather than an arbitrary chamber \(\gamma\),
sets local \(M=V/a_f\) and the Mach angle.

\current \code{FrozenIdealGasTable} already integrates enthalpy and
standard entropy from its bounded \(c_p(T)\) representation.
\code{solve_frozen_nozzle_expansion} closes a sonic throat and Q1D area
stations. Its result explicitly excludes MOC/Rao characteristics and
equilibrium chemistry. Reuse of this module should preserve its data
identity, composition/freeze basis, domain checks and conservation tests.

\proposal Put the property calculation behind a state contract
\[
\mathcal{G}:\quad
(p,T,Y;\hbox{gas mode, data identity})
\longmapsto
(\rho,h,s,a^2,c_p,\hbox{required partials},\hbox{validity}).
\]
Each requested derivative must use the same composition and equilibrium
constraints as the state. Store units, reference enthalpy convention,
species set, property source/version/hash and admissible domain with
the model. Include viscosity and conductivity through a separately
identified transport closure; an EOS alone does not provide them.

For smooth optimization, interpolation regularity also matters. The
present piecewise-linear \(c_p\) law gives consistent integrals, but
derivatives can change at property-grid knots. Any smoother replacement
must preserve \(h_T=c_p\), \(s_T|_p=c_p/T\), positivity and
\(c_p-c_v=R\); independently fitting inconsistent \(h,s,c_p\) surfaces
would defeat the purpose.

\subsection{Generalize the compatibility variables, not only the coefficient}
For smooth isentropic turning, the velocity differential
\begin{equation}
\dd\mathcal{N}=\sqrt{M^2-1}\,\frac{\dd V}{V}
\label{eq:turning}
\end{equation}
defines a path-dependent turning increment before the calorically perfect
specialization. Anderson's Eq.~4.35 gives the corresponding turning
differential; Eqs.~4.39--4.44 specialize and integrate
it for the closed Prandtl--Meyer function
\cite[\S4.14, PDF pp.~186--187]{anderson}.

In the smooth, steady, adiabatic, irrotational/homentropic axisymmetric
case, the velocity-form compatibility retains the geometric source:
\begin{align}
C^+:\quad \dd\theta-\dd\mathcal{N}
&=-\frac{\sin\theta}{Mr}\,\dd\ell,
&\frac{\dd r}{\dd x}&=\tan(\theta+\mu),\\
C^-:\quad \dd\theta+\dd\mathcal{N}
&=+\frac{\sin\theta}{Mr}\,\dd\ell,
&\frac{\dd r}{\dd x}&=\tan(\theta-\mu).
\end{align}
Here \(\ell\) is characteristic arclength oriented downstream, distinct
from thermodynamic entropy \(s\), and \(\mu=\sin^{-1}(1/M)\).
The displayed source form applies for \(r>0\); the axis requires a
separately regularized symmetry limit. Backward reconstruction must
retain the same signed orientation convention.
The original axisymmetric derivation uses \(dp=-\rho V\,dV\) and the
isentropic acoustic derivative before introducing the perfect-gas
Prandtl--Meyer function
\cite[\S11.4, PDF p.~408; \S11.8, PDF pp.~420--422,
especially Eq.~11.51]{anderson}.

For the first frozen-property target, recover local \(T,p,\rho,a\) from
\(h_0=h(T)+V^2/2\), the chosen isentrope and EOS. A turning integral or
its differential then depends on that thermodynamic path. It is not the
old universal \(\nu(M;\gamma)\) with a varying argument inserted into it.
Across streamline entropy differences, finite-rate chemistry,
independently varying composition, heat addition or shocks, additional
transport and compatibility terms are required. Smooth homentropic
equilibrium flow can still use a consistently closed EOS in this framework.
Equations above do not claim to supply that larger system.

\subsection{Starting data, sonic closure and conservation}
The transonic starting line must be consistent with the chosen gas law
and throat curvature. A variable-property downstream solver started from
an unverified constant-gas KL expansion is a mixed formulation. The
original KL coefficient table remains unavailable in the prior audit;
this is a separate source gap to resolve or avoid with a verified
alternative initializer.

The first milestone should therefore include the Q1D sonic state,
throat mass flux, a source-grounded transonic/near-sonic initializer and
convergence away from any numerical supersonic offset. Track the reference
state used to define \(c^*\), Mach and mass flow. Require independent
conservation checks rather than rescaling the starting line to an ideal
one-dimensional mass flow by construction.

For nonequilibrium flow, the sonic location is not in general guaranteed
to coincide with the geometric throat. The original text also makes
clear that finite-rate irreversibility breaks isentropy in an otherwise
inviscid adiabatic nozzle \cite[\S17.11, PDF p.~704]{anderson}. The
current frozen sonic-root assumptions cannot be carried into that
branch without a new derivation.

\subsection{Equilibrium and finite-rate branches}
At equilibrium, composition derivatives contribute to heat capacities:
\begin{equation}
c_{p,\mathrm{eq}}
=\sum_iY_i c_{p,i}
+\sum_i h_i\left(\frac{\partial Y_i}{\partial T}\right)_p .
\end{equation}
The source separately defines frozen and equilibrium heat capacities
\cite[PDF p.~678, Eqs.~17.20--17.25]{anderson}.
Formation-energy conventions matter when \(Y_i\) changes. A chamber
property table in \((P_c,\mathrm{O/F})\) is not a complete
local equilibrium EOS in the nozzle.

For an equilibrium EOS expressed in \((\rho,T)\), the following total
differential is a useful proposed implementation check:
\begin{equation}
a_{\mathrm{eq}}^2
=p_\rho|_T-p_T|_\rho\frac{s_\rho|_T}{s_T|_\rho}.
\end{equation}
This follows by imposing \(\dd s=0\) on a differentiable single-phase
equilibrium branch with fixed elemental inventory and \(s_T|_\rho\ne0\).
It is an identity to test against the thermodynamic backend,
not a replacement for computing consistent equilibrium derivatives.
Anderson's more explicit sound-speed expression includes the
non-perfect-gas derivative corrections \cite[PDF p.~682, Eq.~17.35]{anderson}.

A single-temperature finite-rate ideal-mixture branch would additionally
solve, under its declared inviscid assumptions,
\begin{equation}
\frac{D Y_i}{Dt}=\frac{\dot\omega_i}{\rho},\qquad
\sum_iY_i=1,\qquad \sum_i\dot\omega_i=0.
\end{equation}
The production units are \(\mathrm{kg\,m^{-3}\,s^{-1}}\).
Species transport is source-grounded
\cite[PDF p.~687, Eq.~17.42]{anderson}. A reacting
characteristic intersection needs streamline history as well as its
two acoustic parents. The local Young thesis constructs that third
streamline foot explicitly \cite[PDF pp.~91--96]{young}.

The absolute mixture enthalpy must include formation contributions
consistently. Chemical energy cannot also be added as a second external
heat-release term if already present in that enthalpy. Finite-rate
species and internal-mode relaxation can introduce stiff source
integration and additional sound/characteristic conventions.

\subsection{A source discrepancy that blocks direct finite-rate transcription}
Independent checks of the original Young thesis found unresolved
differences between its earlier EOS/absolute-enthalpy definitions and
the printed chemical-source expression. PDF pp.~87--89 define
\(R=\sum_iY_iR_i\), mass production and \(h_i+h_i^\circ\), whereas
Eqs.~4.22--4.23 print a common mixture \(R\) and
\(h_i-h_i^\circ\). A separate units check found a discrepancy between
Eq.~4.49 and the subsequent coefficient forms on PDF pp.~93--94.

These are recorded as source inconsistencies and independently derived
candidates in \code{source_discrepancies.json}, not as an established
published erratum. Young remains useful evidence for the coupled
characteristic/streamline architecture. Its printed source coefficients
are not implementation-ready until rederived and corroborated against
the original governing formulation. This blocks that transcription,
not the earlier fixed-composition thermally perfect milestone.
